\textbf{Phần 1. Trắc nghiệm lựa chọn}
\setcounterpageref{ex}{0}

\begin{ex} %[3H5N31]
Gam là đơn vị dùng để đo
\choice
{Độ dài}
{Dung tích}
{Diện tích}
{\True Khối lượng}
\loigiai{
Gam là đơn vị đo khối lượng, viết tắt là $\text{g}$.\\
\textbf{Đáp án của câu hỏi là D}
}
\end{ex}

\begin{ex} %[3H5N31]
$1\text{ kg}$ bằng bao nhiêu gam?
\choice
{\True $1000\text{ g}$}
{$10\text{ g}$}
{$100\text{ g}$}
{$10000\text{ g}$}
\loigiai{
$1\text{ kg} = 1000\text{ g}$.\\
\textbf{Đáp án của câu hỏi là A}
}
\end{ex}

\begin{ex} %[3H5H31]
Tính: $350\text{ g} + 250\text{ g} = \dots$
\choice
{$500\text{ g}$}
{$700\text{ g}$}
{$650\text{ g}$}
{\True $600\text{ g}$}
\loigiai{
$350 + 250 = 600\text{ g}$.\\
\textbf{Đáp án của câu hỏi là D}
}
\end{ex}

\begin{ex} %[3H5H31]
Tính: $1000\text{ g} - 450\text{ g} = \dots$
\choice
{$500\text{ g}$}
{\True $550\text{ g}$}
{$650\text{ g}$}
{$450\text{ g}$}
\loigiai{
$1000 - 450 = 550\text{ g}$.\\
\textbf{Đáp án của câu hỏi là B}
}
\end{ex}

\begin{ex} %[3H5H31]
Một gói bánh nặng $200\text{ g}$. Hỏi 4 gói bánh như thế nặng bao nhiêu gam?
\choice
{$600\text{ g}$}
{$1000\text{ g}$}
{$700\text{ g}$}
{\True $800\text{ g}$}
\loigiai{
$200 \times 4 = 800\text{ g}$.\\
\textbf{Đáp án của câu hỏi là D}
}
\end{ex}

\begin{ex} %[3H5H31]
Đổi: $1\text{ kg } 500\text{ g} = \dots\text{ g}$.
\choice
{$1050\text{ g}$}
{\True $1500\text{ g}$}
{$150\text{ g}$}
{$1550\text{ g}$}
\loigiai{
$1000 + 500 = 1500\text{ g}$.\\
\textbf{Đáp án của câu hỏi là B}
}
\end{ex}

\begin{ex} %[3H5V31]
Một quả táo nặng $250\text{ g}$, quả lê nặng hơn quả táo $70\text{ g}$. Cả hai quả nặng bao nhiêu gam?
\choice
{\True $570\text{ g}$}
{$550\text{ g}$}
{$600\text{ g}$}
{$500\text{ g}$}
\loigiai{
Quả lê: $250 + 70 = 320\text{ g}$.\\
Cả hai: $250 + 320 = 570\text{ g}$.\\
\textbf{Đáp án của câu hỏi là A}
}
\end{ex}

\begin{ex} %[3H5V31]
Một túi muối nặng $1\text{ kg}$. Người ta đã dùng hết $350\text{ g}$. Lượng muối còn lại trong túi là
\choice
{$600\text{ g}$}
{$700\text{ g}$}
{\True $650\text{ g}$}
{$750\text{ g}$}
\loigiai{
$1000 - 350 = 650\text{ g}$.\\
\textbf{Đáp án của câu hỏi là C}
}
\end{ex}

\begin{ex} %[3H5V31]
Một hộp trà có 5 gói nhỏ, mỗi gói nặng $20\text{ g}$. Hộp trà đó nặng bao nhiêu gam?
\choice
{$80\text{ g}$}
{\True $100\text{ g}$}
{$120\text{ g}$}
{$90\text{ g}$}
\loigiai{
$20 \times 5 = 100\text{ g}$.\\
\textbf{Đáp án của câu hỏi là B}
}
\end{ex}

\begin{ex} %[3H5V31]
Có $600\text{ g}$ bột mì chia đều vào 6 túi. Mỗi túi đựng số gam bột mì là
\choice
{$80\text{ g}$}
{$120\text{ g}$}
{\True $100\text{ g}$}
{$90\text{ g}$}
\loigiai{
$600 : 6 = 100\text{ g}$.\\
\textbf{Đáp án của câu hỏi là C}
}
\end{ex}

\begin{ex} %[3H5V31]
Tính: $125\text{ g} \times 4 = \dots$
\choice
{\True $500\text{ g}$}
{$450\text{ g}$}
{$550\text{ g}$}
{$600\text{ g}$}
\loigiai{
$125 \times 4 = 500\text{ g}$.\\
\textbf{Đáp án của câu hỏi là A}
}
\end{ex}

\begin{ex} %[3H5C31]
Có 3 quả cam có khối lượng bằng nhau và 1 quả xoài nặng $400\text{ g}$. Biết cả 3 quả cam và quả xoài nặng đúng $1\text{ kg}$. Mỗi quả cam nặng bao nhiêu gam?
\choice
{$180\text{ g}$}
{$220\text{ g}$}
{\True $200\text{ g}$}
{$250\text{ g}$}
\loigiai{
3 quả cam nặng: $1000 - 400 = 600\text{ g}$.\\
Mỗi quả cam: $600 : 3 = 200\text{ g}$.\\
\textbf{Đáp án của câu hỏi là C}
}
\end{ex}

\textbf{Phần 2. Trắc nghiệm đúng sai}
\setcounterpageref{ex}{0}

\begin{ex} %[3H5H31]
Xét tính đúng sai về đơn vị đo gam:
\choiceTF
{\True $1\text{ kg} = 1000\text{ g}$.}
{\True $500\text{ g} + 500\text{ g} = 1\text{ kg}$.}
{$1\text{ kg}$ bông nhẹ hơn $1\text{ kg}$ sắt.}
{\True $1\text{ kg}$ bông có khối lượng bằng $1\text{ kg}$ sắt (đều bằng $1000\text{ g}$).}
\loigiai{
\begin{itemize}
    \item $1\text{ kg} = 1000\text{ g}$: Đúng.
    \item $500 + 500 = 1000\text{ g} = 1\text{ kg}$: Đúng.
    \item Đều là $1\text{ kg}$ nên khối lượng bằng nhau. Ý 3 Sai, Ý 4 Đúng.
\end{itemize}
\begin{itemize}
    \item a) \textbf{Đúng}.
    \item b) \textbf{Đúng}.
    \item c) \textbf{Sai}.
    \item d) \textbf{Đúng}.
\end{itemize}
}
\end{ex}

\begin{ex} %[3H5V31]
Mẹ mua $1\text{ kg}$ thịt heo, chia làm 2 phần: phần làm sườn xào chua ngọt nặng $600\text{ g}$, phần còn lại luộc. Xét tính đúng sai:
\choiceTF
{\True Phần thịt luộc nặng $400\text{ g}$.}
{\True Phần thịt làm sườn xào nặng hơn phần thịt luộc $200\text{ g}$.}
{Phần thịt luộc nặng gấp đôi phần thịt xào.}
{\True Phần thịt làm sườn xào nặng gấp $1{,}5$ lần phần thịt luộc.}
\loigiai{
\begin{itemize}
    \item Thịt luộc: $1000 - 600 = 400\text{ g}$. Đúng.
    \item Hiệu: $600 - 400 = 200\text{ g}$. Đúng.
    \item Thịt luộc ($400\text{ g}$) nhẹ hơn thịt xào ($600\text{ g}$). Ý 3 Sai, Ý 4 Đúng.
\end{itemize}
\begin{itemize}
    \item a) \textbf{Đúng}.
    \item b) \textbf{Đúng}.
    \item c) \textbf{Sai}.
    \item d) \textbf{Đúng}.
\end{itemize}
}
\end{ex}

\textbf{Phần 3. Trả lời ngắn}
\setcounterpageref{ex}{0}

\begin{ex} %[3H5H31]
Tính: $450\text{ g} + 350\text{ g}$.
\shortans{800}
\loigiai{
\textbf{Đáp số:} $800$\\
$450 + 350 = 800\text{ g}$.
}
\end{ex}

\begin{ex} %[3H5H31]
Tính: $1000\text{ g} - 280\text{ g}$.
\shortans{720}
\loigiai{
\textbf{Đáp số:} $720$\\
$1000 - 280 = 720\text{ g}$.
}
\end{ex}

\begin{ex} %[3H5V31]
Có 4 gói kẹo, mỗi gói nặng $150\text{ g}$. Tính tổng khối lượng 4 gói kẹo theo gam.
\shortans{600}
\loigiai{
\textbf{Đáp số:} $600$\\
$150 \times 4 = 600\text{ g}$.
}
\end{ex}

\begin{ex} %[3H5V31]
Một quả đu đủ nặng $800\text{ g}$, quả cam nhẹ hơn quả đu đủ $650\text{ g}$. Hỏi quả cam nặng bao nhiêu gam?
\shortans{150}
\loigiai{
\textbf{Đáp số:} $150$\\
$800 - 650 = 150\text{ g}$.
}
\end{ex}

\begin{ex} %[3H5V31]
Đổi: $1\text{ kg } 250\text{ g}$ ra gam.
\shortans{1250}
\loigiai{
\textbf{Đáp số:} $1250$\\
$1000 + 250 = 1250\text{ g}$.
}
\end{ex}

\begin{ex} %[3H5C31]
Một hộp bánh nặng $500\text{ g}$, riêng vỏ hộp nặng $50\text{ g}$. Trong hộp có 9 chiếc bánh bằng nhau. Hỏi mỗi chiếc bánh nặng bao nhiêu gam?
\shortans{50}
\loigiai{
\textbf{Đáp số:} $50$\\
Bánh nặng: $500 - 50 = 450\text{ g}$. Mỗi chiếc bánh: $450 : 9 = 50\text{ g}$.
}
\end{ex}
